\section{Introducción}
\label{sec:intro}

La resonancia magnética (RM) es el estudio de referencia para evaluar ligamentos, meniscos,
cartílago y hueso de la rodilla. Su lectura exige radiólogos especializados y presenta
variabilidad entre observadores, lo que motiva sistemas de apoyo basados en aprendizaje
profundo~\citep{bien2018mrnet}. Este trabajo aborda la clasificación multietiqueta de doce
hallazgos (lesión del LCA y del LCM, daño del menisco medial y lateral, artrosis medial,
lateral y patelofemoral, derrame, sinovitis, quiste de Baker, contusión ósea y fractura) en
el conjunto RSNA Knee Abnormality Detection.

La dificultad central es la supervisión. Solo 58 de los 4\,407 estudios tienen etiquetas
doradas anotadas por radiólogos; el resto solo trae el reporte en texto libre, escrito en
nueve idiomas por 16 centros de cinco continentes. Las etiquetas de entrenamiento deben, por
tanto, extraerse del texto, con el ruido que eso implica. Siguiendo la reducción de escala del
enunciado, trabajamos con una serie sagital con supresión de grasa por estudio, entrada 2.5D y
una muestra estratificada de 600 estudios, y reservamos los 58 dorados exclusivamente para
la evaluación final.

Nuestras contribuciones son:
\begin{itemize}
  \item un extractor de etiquetas multilingüe basado en reglas, con negación, incertidumbre
        y confianza, validado contra los dorados: exactitud 0,81, es decir, ruido del 19\,\%
        (Sección~\ref{sec:datos});
  \item dos arquitecturas 2.5D con agregación por atención por hallazgo: una CNN diseñada
        desde cero y una ResNet-50 preentrenada; la segunda las supera con una diferencia
        significativa en la evaluación pareada (Secciones~\ref{sec:modelos}
        y~\ref{sec:resultados});
  \item una evaluación con intervalos de confianza, comparación pareada, tres semillas y
        ablaciones de pérdida, manejo del ruido, plano y profundidad reentrenada;
  \item un análisis con Grad-CAM que muestra dónde el modelo localiza la anatomía y dónde
        usa atajos por correlación (Sección~\ref{sec:gradcam}), y uno de sesgo por idioma
        y centro (Sección~\ref{sec:sesgo}).
\end{itemize}
El código es reproducible de extremo a extremo con un solo comando (\texttt{run\_all.py}).
